\section{半阶二维矩形的精确测试上确界}\label{rect:section}

全体首项的面模型允许不同长宽的并列构型。下面计算这些构型在有限阈值下的
真实光晕，说明固定正方形的高阶系数不能直接当作全体最优系数。
本节仍以中心轴平行正方形为算子基，输入则允许是矩形。

\begin{theorem}[矩形类的精确公式]\label{rect:exact}
令 $n=2$、$\beta=1$、$1/\sqrt2<\alpha<1$。
归一化矩形输入为
$E_w=[-1/(2w),1/(2w)]\times[-w/2,w/2]$，$0<w\le1$，故 $|E_w|=1$。
若 $w\le\alpha$，则严格光晕为空；若 $\alpha<w\le1$，置
$s=(w+w^{-1})/2$，则
\[
 |H_{1,\mathcal Q_c}(\alpha,E_w)|=2(s-\alpha)(\alpha^{-1}-s).
\]
因此在所有正有限面积的轴平行矩形上，尺度不变光晕比值的上确界恰为
\begin{equation}\label{rect:sup}
 R(\alpha)=\frac12(\alpha^{-1}-\alpha)^2.
\end{equation}
该上确界不由任一矩形达到；由 $w\downarrow\alpha$ 的矩形族逼近。
\end{theorem}
\begin{proof}
写 $W=1/w\ge w$。任意中心处，边长 $\ell$ 的正方形交叠至多为
$\min(W,\ell)\min(w,\ell)$，除以 $\ell$ 后至多为 $w$，且原点处达到。
这先证明 $w\le\alpha$ 时严格光晕为空。
以下设 $w>\alpha$。中心在矩形外时，至少一个方向的交长至多为 $\ell/2$，
另一个方向至多为 $W$，故归一化平均至多为 $W/2<1/(2\alpha)<\alpha$。
所以光晕全部位于矩形内部。

由反射对称，只需考虑 $0\le x<W/2$、$0\le y<w/2$。
记内接触尺度 $p_1=W-2x$、$p_2=w-2y$，外接触尺度
$r_1=W+2x$、$r_2=w+2y$。
各方向交长随 $\ell$ 依次为 $\ell$、$(p_j+\ell)/2$ 和完整边长。
在接触之间，若两方向均部分相交，平均为
$\ell/4+(p_1+p_2)/4+p_1p_2/(4\ell)$，是凸函数；
其余组合的平均为线性函数、常数或常数加非负常数除以 $\ell$。
故最大值出现在这四个接触尺度中。

令 $d=(W-w)/2$。若 $x+y\le d$，存在尺度既完整包含短边又留在长边内，
平均等于 $w>\alpha$；反过来，这种尺度存在要求 $w+2y\le W-2x$，故平台条件也是必要的。若 $x+y>d$，两个内接触均早于第一个外接触。
在首个内接触处平均为 $\min(p_1,p_2)$，最后内接触处为
$S=(p_1+p_2)/2=s-x-y$，且 $\min(p_1,p_2)\le S$。第一个外接触之后平均不增：
一方向完整捕获、另一方向部分相交时，平均为常数加非负常数除以 $\ell$；
两方向都完整捕获后为 $1/\ell$。
若第一个完整捕获的方向边长为 $b$、中心坐标为 $z$，
该外接触值为 $b(S+2z)/(b+2z)$，是 $S$ 和 $b$ 的凸组合。
由于 $b\ge w>\alpha$，只要某个内接触值超过 $\alpha$，该外接触也超过。
因此只需检测第一个外接触。

置 $A=s-\alpha$，$m_1=2\alpha/W-1$、$m_2=2\alpha/w-1$。
两者均属于 $(0,1)$。当 $y-x\le d$ 时先完整捕获短边，外接触条件等价于
$x+m_2y<A$；当 $y-x\ge d$ 时则等价于 $m_1x+y<A$。
平台区 $x+y\le d$ 自动满足第一个不等式，因为 $A-d=w-\alpha>0$。
两条直线在
\[
 X=(\alpha^{-1}-W)/2>0,\qquad Y=(\alpha^{-1}-w)/2>0
\]
相交，且 $Y-X=d$。此外 $0<X<A<W/2$、$0<Y<A<w/2$。
由上述分区，象限内光晕的闭包是顶点
$(0,0),(A,0),(X,Y),(0,A)$ 的四边形。
其面积为 $A(X+Y)/2$；乘以四，得到定理中的面积式。
边界线段面积为零，故严格水平不改变此面积。

最后 $s\in[1,s_*)$，其中 $s_*=(\alpha+\alpha^{-1})/2$，且
\begin{equation}\label{rect:deficit}
 2(s-\alpha)(\alpha^{-1}-s)=R(\alpha)-2(s_*-s)^2.
\end{equation}
这给出严格递增、上确界及不达到的结论。一般矩形通过平移、统一伸缩和坐标交换
化为上述情形，光晕比值保持不变。
\end{proof}

\begin{corollary}[更强的全体下界与任意阶测试展开]\label{rect:orders}
在上述参数范围内，
\[
 C_{1,\mathcal Q_c}(\alpha)\ge R(\alpha)
 >|H_{1,\mathcal Q_c}(\alpha,Q_0)|,
 \qquad
 R(\alpha)-|H_{1,\mathcal Q_c}(\alpha,Q_0)|
 =\frac{(1-\alpha)^4}{2\alpha^2}.
\]
令 $t=1-\alpha$，任意整数 $N\ge2$ 均有精确余项式
\begin{equation}\label{rect:all-orders}
 R(1-t)=2t^2+\sum_{j=3}^N\frac{j+1}{2}t^j
 +\frac{t^{N+1}[(N+2)-(N+1)t]}{2(1-t)^2}.
\end{equation}
所以此测试上确界的截断误差阶 $t^{N+1}$ 不可提高。
\end{corollary}
\begin{proof}
全体输入包含全部矩形。正方形面积由定理~\ref{rect:exact}取 $w=1$ 得到，
相减即为所示正数。展开
$R(1-t)=\tfrac12[(1-t)^{-2}-2+(1-t)^2]$，求几何级数尾和即得余项式。
\end{proof}

\begin{remark}[可达到的测试序列与量词]\label{rect:scope}
式~\eqref{rect:deficit}还给出
\[
 R(\alpha)-|H(\alpha,E_w)|
 =\frac{(w-\alpha)^2(1-\alpha w)^2}{2\alpha^2w^2}.
\]
例如取 $w=1-t+t^m$、整数 $m\ge2$，误差为 $O(t^{2m+2})$，
因而可由实际严格矩形输入实现任意指定有限阶的上述测试展开。
但临界输入 $w=\alpha$ 的光晕为空，不能交换输入极限和严格水平集。
本节并未证明 $C=R$，也未确定全体三阶系数；因此不能直接将 $5/2$ 称作
全体四阶系数。已经严格排除的是：固定正方形在这些有限阈值下达到全体上确界。
\end{remark}

\subsection{有限矩形并的精确尺度搜索}
\begin{proposition}[二维半阶的有限尺度公式]\label{rect:union-scales}
设 $E$ 是有限个有界轴平行矩形
$R_i=[a_i,b_i]\times[c_i,d_i]$ 的并，各矩形内部互不相交，且 $|E|>0$。
给定中心 $z=(z_1,z_2)$，令 $\Lambda_z$ 为下列数中全部正数的集合：
\[
 2|z_1-a_i|,\quad2|z_1-b_i|,\quad2|z_2-c_i|,\quad2|z_2-d_i|.
\]
则
\[
 M_{1,\mathcal Q_c}\mathbf1_E(z)
 =\max_{\ell\in\Lambda_z}\frac{|E\cap(z+[-\ell/2,\ell/2]^2)|}{\ell}.
\]
中心可以位于 $E$ 外部；不要求输入凸或连通。
\end{proposition}
\begin{proof}
在相邻接触尺度之间，每个坐标的交长都是 $\ell$ 的仿射函数，其斜率属于
$\{0,1/2,1\}$。矩形内部不交，故总交叠面积可直接求和，写成
$A\ell^2+B\ell+C$，其中 $A\ge0$，而 $B,C$ 允许为任意实数。除以 $\ell$ 后是
$F(\ell)=A\ell+B+C/\ell$。若 $C>0$，则 $F$ 凸；若 $C\le0$，则
$F'=A-C/\ell^2\ge0$。因此任一有限区间内的最大值都在端点。
平均量在 $\ell\downarrow0$ 和 $\ell\to\infty$ 时均趋于零，
且 $E$ 有正面积，故某个尺度值为正；两个无界端区间也不能提供比接触尺度更大的值。
交叠连续，所以得到所述有限最大值。矩形边界的零面积交叠不影响求和。
\end{proof}

该公式提供尺度优化的精确有限化，但浮点水平集求根和面积积分仍需要独立的误差控制。
有重叠的矩形表示必须先分割为内部不交的表示，不能直接重复累加面积。
